\begin{proof}

\begin{enumerate}[label={$\langle1\rangle$\arabic*.},leftmargin=*]
\item Let $D$ be the conductor. P4.12 makes $K_T+D$ torsion;
pushing down the point blowups gives $K_{T_0}+D_0$ torsion,
with $D_0$ effective and nonzero. This is a numerically
anticanonical divisor in the sense of Dloussky,
Definition~2.20 of \cite{Dloussky2006}. His Theorem~0.2
therefore gives a global spherical shell. Proposition~2.27
and the square-zero-divisor paragraph after Lemma~2.22
leave the parabolic Inoue or Inoue--Hirzebruch cases,
or a connected rational cycle with attached rational
trees. The latter configuration and $\pi_1(T_0)=\mathbb Z$
are stated explicitly in \cite[\S1, pp.~284--285]{DOT2003}.
We verify the consequences needed
for our conductor.

\item For an Inoue--Hirzebruch surface, the reduced cycle or
union of two cycles is numerically anticanonical. The
curve intersection matrix is negative definite. Subtracting
that divisor from $D_0$ and intersecting with every curve
therefore makes every coefficient difference zero.
Thus $D_0$ is reduced nodal. In the parabolic Inoue case
there is an elliptic curve $E$ and a disjoint rational
cycle $\Gamma$, with
\[
 K_{T_0}+E+\Gamma\sim0,\qquad \Gamma\equiv0.
\]
These relations are stated after Dloussky's Lemma~2.22.
Intersections with $E$ and the rational curves give
$D_0=E+k\Gamma$: the cycle matrix has kernel generated
by the sum of its components. Torsion of $K_{T_0}+D_0$
then makes $\mathcal O((k-1)\Gamma)$ torsion. A nonzero
effective divisor cannot have torsion associated line
bundle on a compact connected reduced complex space:
a power of its defining section would be a nonzero
constant and could have no zero. Hence $k=1$ and
$D_0$ is again reduced nodal.

\item The blowup coefficient $\mu-1$ preserves the reduced
nodal property, as in P4.21. Thus $S$ has ordinary
double curves outside finitely many points. P4.20
makes it normal crossing everywhere. Its normalization
is then immersive and P4.6 gives $e(T)=K_T^2$.
For a class-VII surface, however,
\[
 e(T)=b_2(T)>0,\qquad K_T^2=-b_2(T),
\]
a contradiction. These two exceptional minimal-model
cases are excluded.

\item In the remaining case the support of the conductor is
a connected rational cycle with attached rational trees.
It is the whole connected curve configuration: Dloussky's
Lemma~2.22(2) puts every curve meeting the anticanonical
support in that support, and Proposition~2.27 gives its
connectedness. Its fundamental group is $\mathbb Z$.
The blowups preserve this group by the quotient argument
of P4.21, and preserve the surface fundamental group.
Consequently
\[
 \pi_1(D_{\mathrm{red}})=\mathbb Z,\qquad
 \pi_1(T)=\mathbb Z.
\]
Every normalized component of
$C=\operatorname{Sing}(S)_{\mathrm{red}}$ is rational,
being the finite nonconstant image of a rational conductor
component. The curve is connected. Write $\pi_1(C)=F_b$
for the free group of its point-identification graph.
The conductor pushout gives
\[
 \pi_1(S)=
 \langle t,c_1,\ldots,c_b\mid t^k=w(c_1,\ldots,c_b)\rangle.
\]
Here $k$ records the conductor-cycle map to $\pi_1(T)$,
and $w$ records its map to $\pi_1(C)$. Neither map is
assumed to be an isomorphism.

\item Virtual abelianness from P4R.1 forces $b\le1$.
If $k\ne0$ and $w\ne1$, both cyclic amalgam maps are
injective; alternating nonidentity coset representatives
give the reduced normal form and show that $F_b$ injects
in the amalgam. If $w=1$, the group is
$(\mathbb Z/k\mathbb Z)*F_b$, with
$\mathbb Z/0\mathbb Z=\mathbb Z$, and $F_b$ again
injects. A nonabelian free subgroup is impossible in a
virtually abelian group. Finally, if $k=0$ and $w\ne1$,
the group is
\[
 \mathbb Z*\bigl(F_b/\langle\!\langle w\rangle\!\rangle\bigr).
\]
A free product of an infinite cyclic group and a nontrivial
group contains a nonabelian free subgroup: for a nontrivial
$g$ in the second factor, $t$ and $gtg^{-1}$ have no
nonempty reduced relation. Thus the second factor must
be trivial. Its abelianization has rank at least $b-1$,
which again gives $b\le1$.

\item For the number $c\ge1$ of components of $C$, we obtain
\[
 e(C)=c+1-b\ge c\ge1.
\]
By P4.13, $e(S_{\mathrm{reg}})\ge0$, so
$e(S)=e(S_{\mathrm{reg}})+e(C)>0$.

\end{enumerate}
\end{proof}
